\documentclass[11pt,a4paper]{article}

% ---------- packages ----------
\usepackage[utf8]{inputenc}
\usepackage[T1]{fontenc}
\usepackage{lmodern}
\usepackage{amsmath,amssymb,amsthm}
\usepackage{mathtools}
\usepackage{physics}
\usepackage{bm}
\usepackage{geometry}
\geometry{margin=2.5cm}
\usepackage{enumitem}
\usepackage{booktabs}
\usepackage{caption}
\usepackage[dvipsnames]{xcolor}
\usepackage{hyperref}
\hypersetup{colorlinks,linkcolor=MidnightBlue,citecolor=OliveGreen,urlcolor=NavyBlue}
\usepackage{cleveref}
\usepackage{tikz}
\usetikzlibrary{calc,decorations.pathreplacing,patterns,arrows.meta}

% ---------- theorem styles ----------
\newtheorem{theorem}{Theorem}
\newtheorem{lemma}[theorem]{Lemma}
\newtheorem{proposition}[theorem]{Proposition}
\newtheorem{corollary}[theorem]{Corollary}
\theoremstyle{definition}
\newtheorem{definition}[theorem]{Definition}
\theoremstyle{remark}
\newtheorem{remark}[theorem]{Remark}
\newtheorem{claim}[theorem]{Claim}

% ---------- shortcuts ----------
\newcommand{\Sab}{S_{\mathrm{ab}}}
\newcommand{\Sinc}{S_{\mathrm{inc}}}
\newcommand{\Tzero}{T_0}
\newcommand{\nhat}{\hat{\vb{n}}}
\newcommand{\khat}{\hat{\vb{k}}}
\newcommand{\ntilde}{\tilde{n}}
\newcommand{\SAR}{\ensuremath{\mathrm{SAR}}}
\newcommand{\psSAR}{\ensuremath{\mathrm{psSAR}_{10\mathrm{g}}}}
\newcommand{\rhom}{\rho_{\mathrm{m}}}
\newcommand{\Lact}{L_{\!*}}

\title{%
  When the Cauchy projection formula for $\psSAR$ fails:\\
  side-entry constraints under the IEC/IEEE 62704 cube rule%
}
\author{Robin Wydaeghe (with Claude Opus 4.7)}
\date{May 2026}

\begin{document}
\maketitle

\begin{abstract}
A previous note (\texttt{psSAR10g.tex}) derived a closed-form thin-skin
expression for the peak spatial-average SAR over the ICNIRP/IEEE 10\,g
cube,
\[
  \langle\SAR\rangle_{\mathrm{cube}}
  = \frac{\Sinc\,\Tzero}{\rhom\,L}\bigl(|k_x|+|k_y|+|k_z|\bigr),
  \qquad L = (m/\rhom)^{1/3}=21.5\ \mathrm{mm},
\]
on the basis of an energy-conservation argument: the absorbed power
inside the cube equals the projected area $A_\perp(\khat) =
L^2(|k_x|+|k_y|+|k_z|)$ of the cube on the wavefront, multiplied by
$\Sinc\Tzero$. We revisit this derivation in the light of the actual
IEC/IEEE 62704-1 averaging-volume rules (axis-aligned cube, mass
fixed at 10\,g, at most 10\% air by volume, no face entirely in air)
and find that the formula is correct only in a narrow regime that does
not include the dominant compliance scenario of plane-wave illumination
on a smooth body. For face-flush axis-aligned cubes on a planar body,
five of the six cube faces are tissue-facing, the side-entry channel
is essentially closed, and the absorbed power inside the cube is
$\Sinc\Tzero L^2\cos\theta_i$ to within layer-induced corrections,
where $\theta_i$ is the local incidence angle on the body surface.
A numerical 62704-style optimisation on the Thelonious phantom across
four frequencies (2.45, 10, 28, 60\,GHz), two polarisations, and five
incidence directions confirms the analysis: the Cauchy formula
overpredicts $\psSAR$ by a median of $3\times$ (and up to $50\times$
at grazing incidence), while the cosine-law replacement
$\langle\SAR\rangle = \Sab/(\rhom L)$ stays within $\sim 50\%$ of
numerical, with the residual bias driven by the angle $\gamma$ between
the body normal and the cube axes (zero for axis-aligned, growing to
$+50\%$ at $\gamma=45^\circ$). We give the corrected formula, identify
the regime in which the Cauchy expression remains accurate (sharp
convex corners or small anatomy that the cube engulfs), and recommend
that the corollary in \texttt{paper.tex} be replaced with the cosine
law (with appropriate geometry caveat) or removed.
\end{abstract}

\tableofcontents


% ======================================================================
\section{The standardized averaging volume}
\label{sec:standard}
% ======================================================================

\subsection{IEC/IEEE 62704-1 cube rule, restated}

The peak spatial-average SAR over 10\,g of tissue is computed
according to IEC/IEEE 62704-1:2017 (FDTD) and IEC/IEEE 62704-4:2020
(FEM), using the averaging-volume construction in IEEE C95.3-2021
Annex~E. Three rules are operationally relevant:

\begin{enumerate}[label=(R\arabic*),leftmargin=2.5em]
  \item\label{R:axis} The cube is \textbf{axis-aligned} with the
    global Cartesian voxel grid. It does not rotate to follow the body
    surface.
  \item\label{R:mass} The cube is \textbf{grown} from a seed voxel
    until the enclosed tissue mass equals the prescribed
    $m=10$\,g (within a small tolerance).
  \item\label{R:air} The cube is \textbf{valid} iff (a) the enclosed
    air fraction does not exceed $10\%$ of the cube volume, and (b)
    no single face is entirely in air. (Plain IEEE C95.3 omits the
    $10\%$ rule and uses only the no-air-face condition; CST
    documents both methods.)
\end{enumerate}

\noindent
The first consequence is that the cube side is not fixed at
$L=(m/\rhom)^{1/3} = 21.54$\,mm. With the maximum 10\% air, the
cube side grows to
\begin{equation}\label{eq:Lact}
  \Lact = \Bigl(\frac{m}{0.9\,\rhom}\Bigr)^{1/3}
  = 22.31\;\mathrm{mm},
\end{equation}
i.e.\ up to $3.6\%$ larger than the all-tissue value. The mass is
$10$\,g by construction, so this is purely a volume effect.

The second consequence is geometric: under~\ref{R:axis}, the cube
sits in a plane wherever the body's local surface normal is not
parallel to a global axis. Most of the body's surface satisfies this
condition, so the body cuts the cube as a tilted slab.

\subsection{Optimal valid placement on a planar body}

Take a half-space body $\{z<0\}$ with horizontal surface and an
axis-aligned cube of side $L$. Two regimes appear depending on
whether the body normal aligns with a cube axis ($\gamma=0$) or not
($\gamma>0$).

\paragraph{Face-flush ($\gamma=0$): no-face-air rule binds.}
For axis-aligned cube on a body whose local normal is parallel to a
cube axis, the cube straddles the surface in the obvious way: top face
above (in air), bottom face below (in tissue). If we try to push the
cube up to allow $f_{\mathrm{air}}>0$, the entire top face moves into
air and rule R3b (no entirely-air face) is violated. The optimum is
therefore the limiting placement with the cube top exactly at the
surface, $f_{\mathrm{air}}=0$ and
\begin{equation}\label{eq:opt-faceflush}
  L = L_{\mathrm{nominal}} = (m/\rhom)^{1/3} = 21.54\;\mathrm{mm}.
\end{equation}

\paragraph{Tilted ($\gamma>0$): 10\%-air rule binds.}
For an axis-aligned cube on a body whose normal makes angle
$\gamma>0$ with every cube axis, the body surface enters the cube as a
tilted plane and the air region inside the cube becomes a wedge that
sits across multiple faces. No single face is entirely in air, so
rule R3b is satisfied even when $f_{\mathrm{air}}\approx 10\%$. The
10\%-air constraint binds, giving
\begin{equation}\label{eq:opt-Lstar}
  L = \Lact = \Bigl(\frac{m}{0.9\,\rhom}\Bigr)^{1/3}
  = 22.31\;\mathrm{mm},
\end{equation}
with cube footprint $\Lact^2 = 4.98$\,cm$^2$.

\medskip\noindent
The face-flush case ($L_{\mathrm{nominal}}$) and the tilted case
($\Lact$) differ in cube side by $\Lact/L_{\mathrm{nominal}} = 1.036$
and in the $1/(\rhom L)$ prefactor of the cosine formula by $0.965$,
so the two limits set tight bounds on which $L$ is operative. The
numerical search in \cref{sec:numerics-thelonious} confirms both
regimes appear on Thelonious: for surface points with
$\gamma\lesssim 5^\circ$ the optimum is the face-flush
\cref{eq:opt-faceflush}; for the rest (the great majority on a smooth
body) the optimum is the 10\%-air \cref{eq:opt-Lstar}.


% ======================================================================
\section{The Cauchy projection corollary}
\label{sec:corollary}
% ======================================================================

The corollary in question, as stated in \texttt{paper.tex}
(Cor.~\texttt{cor:cube} therein) and in \texttt{psSAR10g.tex} (eq.~52), is
\begin{equation}\label{eq:cauchy-form}
  \langle\SAR\rangle_{\mathrm{cube}}
  \stackrel{?}{=} \frac{\Sinc\,\Tzero}{\rhom\,L}\,
    \bigl(|k_x|+|k_y|+|k_z|\bigr),
  \qquad L = (m/\rhom)^{1/3}.
\end{equation}
The argument is energy conservation in the thin-skin limit: every
photon entering the cube's air-facing faces is absorbed inside the
cube, so $P_{\mathrm{abs}}=\Sinc\Tzero A_\perp(\khat)$, with
$A_\perp$ the projected area of the cube on the wavefront. For an
axis-aligned cube,
\begin{equation}\label{eq:cauchy-area}
  A_\perp(\khat) = L^2\bigl(|k_x|+|k_y|+|k_z|\bigr).
\end{equation}
Dividing by $m=\rhom L^3$ gives \cref{eq:cauchy-form}.

The implicit premise is that the body fills the cube's cross-section
perpendicular to $\khat$, so that every photon entering the cube's
geometric shadow is intercepted by tissue. The standalone note
acknowledges this as a caveat for ``small or protruding anatomy''. The
remainder of this report shows that the premise also fails for the
\emph{opposite} extreme --- a smooth body large compared to the cube
--- which is precisely the regime where the formula is most often
invoked (torso, limb, head). The resulting overprediction is
substantial.


% ======================================================================
\section{A face-by-face flux audit}
\label{sec:audit}
% ======================================================================

Place the 10\%-air valid cube of \cref{eq:opt-Lstar} on a horizontal
half-space body. Take a plane-wave incidence direction
$\khat=-(\sin\theta,0,\cos\theta)$ with $\theta\in[0,\pi/2)$ the angle
from the cube's vertical axis. The cube is
$[0,\Lact]\times[0,\Lact]\times[-0.9\Lact, 0.1\Lact]$ and the body
surface is the plane $z=0$ inside the cube footprint.

\subsection{Air-facing vs tissue-facing}

The cube has six faces. Their character under \cref{eq:opt-Lstar}
is:

\begin{center}
\begin{tabular}{@{}llll@{}}
\toprule
Face & Outward normal & Region of face above $z=0$ & Character \\
\midrule
Top    & $+\hat{\vb z}$ & Entire face       & Air-facing \\
Bottom & $-\hat{\vb z}$ & None              & Tissue-facing \\
$x=0$  & $-\hat{\vb x}$ & Strip $z\in[0,0.1\Lact]$ & Mixed \\
$x=\Lact$ & $+\hat{\vb x}$ & Strip $z\in[0,0.1\Lact]$ & Mixed \\
$y=0$  & $-\hat{\vb y}$ & Strip $z\in[0,0.1\Lact]$ & Mixed \\
$y=\Lact$ & $+\hat{\vb y}$ & Strip $z\in[0,0.1\Lact]$ & Mixed \\
\bottomrule
\end{tabular}
\end{center}

The four side faces are split: a thin $\Lact\times 0.1\Lact$ strip at
the top is in air; the lower $\Lact\times 0.9\Lact$ region is in
tissue. Five of the six faces have most or all of their area immersed
in tissue.

\subsection{Side entry through tissue strips is closed}

The crucial point: the tissue strips of the side faces do not channel
new wave power into the cube. Consider a photon present in tissue at
the right-side tissue strip $\{x=\Lact,\,z<0\}$. It must have crossed
the body surface $z=0$ at some location, then propagated in tissue to
reach the cube boundary. In the thin-skin regime $\alpha L\gg 1$ the
characteristic path length in tissue is $1/(2\alpha)\approx 0.5$\,mm
(skin, 28\,GHz), much smaller than the cube side $L=21.5$\,mm. By
Snell refraction into a high-index medium ($|\ntilde|\gtrsim 4$ above
6\,GHz) the in-tissue propagation angle from normal is
$\theta_t\approx\theta/|\ntilde|$, so a photon entering tissue at
$(\Lact+\delta,y,0)$ (just to the right of the cube) travels nearly
straight down and reaches $x=\Lact$ at depth
$z\approx -\delta\cot\theta_t \approx -\delta\,|\ntilde|/\sin\theta$,
with amplitude attenuated by $e^{-\alpha\delta\,\cot\theta_t}$.
Integrated over $\delta\geq 0$, the inward Poynting flux through
the tissue strip of the right face is
\begin{equation}\label{eq:side-channel-bound}
  \Phi^{\mathrm{tissue\,strip}}_{x=\Lact}
  \;\lesssim\; \Sinc\,\Tzero\,\sin\theta\;\Lact\,
    \frac{\sin\theta_t}{2\alpha}
  \;=\; \frac{\Sinc\,\Tzero\,\Lact}{2\alpha\,|\ntilde|}\,\sin^2\theta,
\end{equation}
which at 28\,GHz is $\sim 10^{-4}\,\Sinc\Tzero\Lact^2$ per face.
Compared to the top-face contribution $\Sinc\Tzero\Lact^2\cos\theta$,
the tissue side-channel adds at most $\sim 0.1\%$ across all
$\theta$. We therefore drop it.

\subsection{Direct accounting via the body surface inside the footprint}

The cleanest way to compute the absorbed power inside the cube is to
note that, in the thin-skin limit, all wave power transmitted through
the body surface inside the cube footprint is absorbed inside the
cube (apart from the negligible side leakage of
\cref{eq:side-channel-bound}). The body surface inside the footprint
has area $\Lact^2$ (a horizontal $\Lact\times\Lact$ square at
$z=0$). The local absorbed power density is
$\Sab(\rr)=\Sinc\Tzero\cos\theta_i$, with $\theta_i=\theta$ here
(planar body, $\hat{\vb n}=+\hat{\vb z}$).

Hence
\begin{equation}\label{eq:correct-Pabs}
  P_{\mathrm{abs}}^{\mathrm{cube}}
  = \int_{\Sigma\cap\mathcal{C}} \Sab\,\diff A
  = \Sinc\,\Tzero\,\cos\theta\;\Lact^2,
\end{equation}
not $\Sinc\Tzero\Lact^2(|\sin\theta|+|\cos\theta|)$ as the Cauchy
formula would predict.

\subsection{Cross-check via inward flux through air-facing faces}

The same answer follows from the Poynting-flux picture used in the
original derivation. The inward Poynting flux through the air-facing
parts of the cube boundary is
\begin{equation}\label{eq:air-flux-correct}
  \Phi^{\mathrm{in}}_{\mathrm{air}}
  = \underbrace{\Sinc\,\Lact^2\cos\theta}_{\text{top face}}
  + \underbrace{\Sinc\,(\Lact\times 0.1\Lact)\sin\theta}_{\text{right air strip}}
  = \Sinc\Lact^2\bigl(\cos\theta + 0.1\sin\theta\bigr).
\end{equation}
Of this inflow, only the fraction that hits the body surface
\emph{inside} the cube footprint contributes to
$P_{\mathrm{abs}}^{\mathrm{cube}}$. A photon entering through the
top face at $(x_t,y,0.1\Lact)$ reaches the body at
$x_b = x_t-0.1\Lact\tan\theta$, which lies inside $[0,\Lact]$ only
for $x_t \in [0.1\Lact\tan\theta,\Lact]$. The fraction is
$1-0.1\tan\theta$. A photon entering the right air strip at
$(\Lact,y,z_e)$ reaches the body at
$x_b=\Lact-z_e\tan\theta\in[\Lact-0.1\Lact\tan\theta,\Lact]$,
all inside the footprint. Thus
\begin{align}
  P_{\mathrm{abs}}^{\mathrm{cube}}
  &= \Tzero\bigl[\Sinc\Lact^2\cos\theta\,(1-0.1\tan\theta)
     + \Sinc\,0.1\Lact^2\sin\theta\bigr] \notag\\
  &= \Sinc\,\Tzero\,\Lact^2\bigl[\cos\theta - 0.1\sin\theta + 0.1\sin\theta\bigr]
  = \Sinc\,\Tzero\,\Lact^2\cos\theta,
  \label{eq:flux-cancellation}
\end{align}
in agreement with \cref{eq:correct-Pabs}. The contribution that the
Cauchy formula assigns to the side faces ($\Sinc\Tzero\Lact^2\sin\theta$
in this geometry) is missing the factor $0.1$ from the actual air-strip
height \emph{and} is partly cancelled by the wastage of top-face flux
that misses the cube footprint. The two effects do not add up to the
$\sin\theta$ term that appears in $|k_x|+|k_z|$.

\begin{remark}[Where the Cauchy bookkeeping breaks]
The Cauchy projection $A_\perp=L^2(|k_x|+|k_y|+|k_z|)$ counts the
full geometric shadow of the cube as if the cube were an isolated
opaque object suspended in air with body filling its shadow on every
forward-facing side. For a 62704-valid cube (mostly inside the body),
five of the six faces are immersed in tissue and intercept no fresh
flux; the inflow comes essentially through the single air-facing
face (top), with negligible corrections from the air strips on the
other faces. The result is $\propto\cos\theta_i$, not
$\propto |k_x|+|k_y|+|k_z|$.
\end{remark}


% ======================================================================
\section{Numerical comparison}
\label{sec:numerics}
% ======================================================================

For a horizontal half-space body and an optimally placed valid cube,
the corrected and Cauchy expressions for $\psSAR/(\Sinc\Tzero)$ are
compared in \cref{tab:compare}. Skin tissue at 28\,GHz, $\rhom=1000$
kg/m$^3$, $\Lact=22.31$\,mm, $L=21.54$\,mm.

\begin{table}[ht]
\centering
\caption{$\psSAR$ on a flat horizontal body for an optimal axis-aligned
cube, plane-wave incidence at angle $\theta$ from the cube vertical.
``Corrected'' uses \cref{eq:correct-Pabs} with $\Lact=22.31$\,mm
(10\% air saturation). ``Cauchy'' uses \cref{eq:cauchy-form} with
$L=21.54$\,mm. Errors are
$(\text{Cauchy}-\text{Corrected})/\text{Corrected}$.}
\label{tab:compare}
\begin{tabular}{rcccc}
\toprule
$\theta$ & $\cos\theta$ & Corrected
  $\;[\Sinc\Tzero\,\mathrm{W/kg}]$
  & Cauchy $[\Sinc\Tzero\,\mathrm{W/kg}]$
  & Error \\
\midrule
$0^\circ$  & 1.000 & 0.0498 & 0.0465 & $-7\%$ \\
$15^\circ$ & 0.966 & 0.0481 & 0.0586 & $+22\%$ \\
$30^\circ$ & 0.866 & 0.0431 & 0.0635 & $+47\%$ \\
$45^\circ$ & 0.707 & 0.0352 & 0.0658 & $+87\%$ \\
$60^\circ$ & 0.500 & 0.0249 & 0.0635 & $+155\%$ \\
$75^\circ$ & 0.259 & 0.0129 & 0.0586 & $+355\%$ \\
\bottomrule
\end{tabular}
\end{table}

The Cauchy formula is biased low by $7\%$ at normal incidence (an
$\Lact$ vs $L$ accounting issue) and biased high by $50\%$ to
factor-of-three at oblique incidence (the side-entry overcounting).
Crucially, the Cauchy formula is non-monotone in $\theta$, peaking
near $\theta=45^\circ$, whereas the actual $\psSAR$ decreases
monotonically with the incidence angle. The two expressions
disagree even on the qualitative direction of the answer.

For a generic $\khat$ with all three components nonzero, the same
analysis gives, with $\theta_i$ the local incidence angle on the
body surface (here $\cos\theta_i = -\nhat\cdot\khat = -k_z$ for
horizontal body) and $\nhat=+\hat{\vb z}$,
\begin{equation}\label{eq:correct-3d}
  \langle\SAR\rangle^{\mathrm{flat}}_{\mathrm{cube}}
  = \frac{\Sinc\,\Tzero}{\rhom\,\Lact}\,
    \pospart{\nhat\cdot(-\khat)}
  = \frac{\Sab(\rr)}{\rhom\,\Lact},
\end{equation}
which is the exact analogue of the surface-SAR formula
$\SAR(\rr,0)=2\alpha\Sab/\rhom$ but with the depth scale set by the
cube extent $\Lact$ rather than the skin depth $1/(2\alpha)$. The
Cauchy projection factor is absent.


% ======================================================================
\section{When the Cauchy formula does hold}
\label{sec:limits}
% ======================================================================

The Cauchy formula's premise --- the body fills the cube shadow on
every forward-facing face --- is satisfied in three regimes.

\subsection{Sharp convex corners}\label{sec:corner}

If the body has a convex corner and the cube fits snugly into it,
two or three forward-facing faces of the cube can simultaneously be
air-facing. For an idealized rectangular corner (body fills octant
$\{x<0,y<0,z<0\}$), the cube
$[0,\Lact]^3$ would have its three positive-axis faces all
air-facing and its three negative-axis faces tissue-facing. Inward
flux through the air-facing faces would equal
$\Sinc\Lact^2(|k_x|+|k_y|+|k_z|)$ for incidence with all components
positive (i.e., from $-\khat$ in the all-positive octant), and the
Cauchy formula would be exact.

This is, however, an unphysical limit. Real anatomy has rounded
features with curvature radius $R\gg\Lact$ (torso, limb) or
$R\sim\Lact$ (fingertip, ear edge). The right-angle corner is an
upper bound that the body never approaches. \emph{An interior corner
on the body fails this scenario in the opposite direction: the cube
cannot extend into the corner because most of its volume would be
in tissue from multiple directions, but air-facing flux remains
single-direction.}

\subsection{Body fully engulfed by the cube}\label{sec:engulf}

If a slender body part (a finger, an ear cartilage sheet) fits
entirely inside the cube, the cube ``shadows'' the part from every
forward-facing direction and the projected-area argument applies. For
a finger of radius $r\approx 8$\,mm and $\Lact\approx 22$\,mm, the
finger fills perhaps $\pi r^2 / \Lact^2\approx 41\%$ of the cube
cross-section, so the formula is approximately right (within the
40\% it omits). For ear cartilage of thickness $1$--$2$\,mm, the
fill fraction is below $10\%$ and the Cauchy formula overpredicts
by an order of magnitude, but in the opposite direction from the
\cref{sec:audit} smooth-body case: the cube simply doesn't contain
much absorbing material. The 62704 air-fraction rule tends to rule
out such cubes anyway, since they violate the $\geq 90\%$ tissue
criterion.

\subsection{Axis-aligned incidence}\label{sec:axisaligned}

If $\khat$ is along a global axis (one of $\pm\hat{\vb x}$,
$\pm\hat{\vb y}$, $\pm\hat{\vb z}$), $A_\perp = L^2$ regardless of
body geometry, and the formula reduces to
$\langle\SAR\rangle = \Sinc\Tzero/(\rhom L)$. For a smooth body with
the cube placed face-flush at the wave-incidence side, this matches
\cref{eq:correct-3d} (with $\theta_i=0$) up to the $\Lact$ vs $L$
factor, which is the $7\%$ underprediction in \cref{tab:compare}.

\medskip\noindent
Outside these three regimes, the Cauchy formula does not describe
the standardized $\psSAR$.


% ======================================================================
\section{Numerical validation on Thelonious}
\label{sec:numerics-thelonious}
% ======================================================================

The analytical argument above predicts a specific failure mode for the
Cauchy formula at oblique incidence on a smooth body, and a $\sim$few-\%
match for the cosine law in face-flush geometry. We now check both
against a numerical 62704-style cube optimisation on the
Thelonious phantom under plane-wave illumination, with a three-layer
planar tissue model (skin\,2\,mm / fat\,10\,mm / muscle$_\infty$,
IT'IS-style properties).

\paragraph{Setup.} We sample 15 surface points on Thelonious. At each
point we treat the body locally as a half-space with the surface normal
taken from the STL face. We sweep four frequencies (2.45, 10, 28,
60\,GHz), two polarisations (TE and TM), and five plane-wave incidence
directions (one straight down, one frontal, one side, one 45$^\circ$
above the horizon, one along the body diagonal). For each
configuration we compute the layered Fresnel transmission via a 3-layer
transfer matrix, the SAR depth profile in each tissue layer, and the
maximum mass-averaged SAR over an axis-aligned cube whose placement is
optimised under all three 62704 rules (cube grown to enclose 10\,g,
$\le 10\%$ air by volume, no entirely-air face). The cube integration
uses Monte Carlo (200\,k samples in the cube, focused scan of cube
centres around the face-flush optimum).
After dropping configurations with grazing or back-side incidence
(280 valid records), the headline numbers are:

\begin{table}[ht]
\centering
\caption{Median ratio of closed-form to numerical $\psSAR$ across the
sweep (280 records, Thelonious).}
\label{tab:thelonious-summary}
\begin{tabular}{rccc}
\toprule
$f$ [GHz] & Cauchy / num & CosL$_*$ / num & CosL$_{\mathrm{nom}}$ / num \\
\midrule
2.45  & 3.12 & 1.33 & 1.38 \\
10    & 2.53 & 1.06 & 1.09 \\
28    & 3.06 & 1.31 & 1.35 \\
60    & 3.19 & 1.37 & 1.41 \\
\midrule
all   & 3.09 & 1.30 & 1.30 \\
\bottomrule
\end{tabular}
\end{table}

\paragraph{Cauchy is dramatically wrong.} The median Cauchy/numerical
ratio is 3$\times$, and at grazing incidence ($\cos\theta_i \to 0$) it
diverges to 10$\times$--50$\times$. The $|k_x|+|k_y|+|k_z|$ factor
predicts the largest psSAR in directions where the cube would intercept
the most flux \emph{as a free-floating object in air}, but those are
exactly the directions where the body surface hardly intersects the
cube at all. The two pictures fight each other and Cauchy loses.

\paragraph{Cosine law is much closer but not exact.} Both variants of
the cosine law (with $L_*$ or with $L_{\mathrm{nominal}}$) sit within
the same envelope: median ratios near 1.0--1.5, no divergence at
grazing incidence. The 10\,GHz column is the closest to unity (median
1.06--1.09), because at 10\,GHz the SAR depth profile fills the
$L=21.5$\,mm cube nicely, so even a tilted depth distribution averages
out close to the nominal. At 28 and 60\,GHz the ratio rises to
1.30--1.40 because in the thin-skin regime a tilted cube's
near-surface samples shift below the SAR peak and the mass average
falls. At 2.45\,GHz the ratio rises again from layered standing-wave
effects below the 6\,GHz boundary (already known to break thin-skin).

\paragraph{The dominant residual is the body-normal vs cube-axis tilt
$\gamma$.} Binning by the angle $\gamma$ between the local body normal
and the nearest cube axis:

\begin{table}[ht]
\centering
\caption{Median ratio CosL$_{\mathrm{nom}}$/numerical, binned by the
tilt $\gamma$ between body normal and nearest cube axis. The cube is
axis-aligned in the global frame; $\gamma$ depends only on local body
geometry.}
\label{tab:gamma-bins}
\begin{tabular}{rccccc}
\toprule
$\gamma$ & 2.45\,GHz & 10\,GHz & 28\,GHz & 60\,GHz & median \\
\midrule
5--15$^\circ$  & 1.16 & 0.94 & 0.95 & 0.95 & 0.96 \\
15--25$^\circ$ & 1.26 & 1.03 & 1.19 & 1.24 & 1.20 \\
25--35$^\circ$ & 1.35 & 1.09 & 1.33 & 1.40 & 1.33 \\
35--45$^\circ$ & 1.43 & 1.14 & 1.39 & 1.42 & 1.40 \\
45--55$^\circ$ & 1.46 & 1.16 & 1.45 & 1.50 & 1.46 \\
\bottomrule
\end{tabular}
\end{table}

The cosine law is essentially exact ($\le 5\%$ error) when
$\gamma \le 15^\circ$, and degrades monotonically with $\gamma$, reaching
$+50\%$ overshoot at $\gamma=50^\circ$. This is a pure geometric
effect: an axis-aligned cube on a tilted body samples a depth
distribution that is biased away from the surface, so the mass average
over the cube falls below the face-flush prediction. On a complex body
like Thelonious the median $\gamma$ is around 25$^\circ$
(few surface points are precisely axis-aligned).

\paragraph{What about the 10\%-air optimum?} The numerical search
finds that the optimal placement saturates the air constraint
(median $f_{\mathrm{air}} = 0.09$, median $L = 22.2$\,mm) for every
geometry where the body normal is not exactly aligned with a cube axis.
For axis-aligned ($\gamma \to 0$) face-flush placements, the no-face-air
rule rules out the 10\%-air configuration and the optimum collapses to
$L = L_{\mathrm{nominal}} = 21.54$\,mm with $f_{\mathrm{air}} = 0$.
The two cosine variants therefore correspond to the limiting cases:
$L_{\mathrm{nominal}}$ for axis-aligned smooth body, $L_*$ for everywhere
else. They differ by $1.07\times$, but the geometric tilt residual of
$1.0$--$1.5\times$ dominates either way.

\paragraph{Visual summary.}
\Cref{fig:scatter} shows the
formula--vs--numerical scatter across all 280 configurations. The
Cauchy panel clearly separates above the diagonal across two orders of
magnitude in $\psSAR$. The two cosine panels collapse onto a tight band
just above the diagonal. \Cref{fig:ratio} shows the
formula/numerical ratio against the incidence cosine and confirms that
Cauchy diverges to $50\times$ at grazing while the cosine law stays in
$[1, 2]$.

\begin{figure}[ht]
\centering
\includegraphics[width=\textwidth]{scripts/psSAR10g_validation_out/scatter_formula_vs_numerical.png}
\caption{Closed-form $\psSAR_{10\mathrm{g}}$ vs numerical 62704
maximum-over-placement, on Thelonious with the three-layer planar
tissue model. Dashed line is $y=x$. Cauchy projection (left) is
2--50$\times$ above the diagonal across two decades. Cosine law with
$L_*$ (centre) and with $L_{\mathrm{nominal}}$ (right) collapse onto a
tight band just above the diagonal. Colour: incidence cosine
$\cos\theta_i$.}
\label{fig:scatter}
\end{figure}

\begin{figure}[ht]
\centering
\includegraphics[width=\textwidth]{scripts/psSAR10g_validation_out/ratio_vs_mu.png}
\caption{Formula/numerical ratio vs incidence cosine. Cauchy diverges
to 50$\times$ at grazing incidence. Cosine law (both variants)
saturates around 1.0--1.5, with the 10\,GHz line (purple/pink) the
closest to unity since the cube depth ($L=21.5$\,mm) is comparable to
the SAR penetration there.}
\label{fig:ratio}
\end{figure}


% ======================================================================
\section{Why this matters and what to do}
\label{sec:implications}
% ======================================================================

\subsection{Implication for the paper}

The corollary in \texttt{paper.tex} is presented as a closed-form
shortcut for the IEC/IEEE 62704 quantity, valid in the thin-skin
limit $\alpha L\gg 1$. The thin-skin limit is correctly invoked, but
the side-entry channel is silently assumed to be open. For the
plane-wave-on-torso scenario that drives compliance assessments, the
formula overstates $\psSAR$ by $50\%$ to a factor of three across
realistic incidence angles (\cref{tab:compare}). This is not a
controlled approximation; it is a structural error in the energy
bookkeeping.

The error is masked at axis-aligned incidence (where the formula
happens to be approximately right by virtue of $|k_x|+|k_y|+|k_z|=1$
matching $\cos\theta_i=1$) and could go unnoticed in a compliance
table dominated by axis-aligned columns. It will reappear as soon as
a base-station antenna is tilted, the user's body rotates, or
multipath delivers oblique components.

\subsection{Recommended replacement}

The corrected closed-form expression for a smooth-body face-flush
configuration is
\begin{equation}\label{eq:replacement}
  \boxed{%
    \langle\SAR\rangle^{\mathrm{cube,smooth}}
    = \frac{\Sab(\rr)}{\rhom\,\Lact}
    = \frac{\Sinc\,\Tzero}{\rhom\,\Lact}\,
      \pospart{\nhat(\rr)\cdot(-\khat)},
  }
\end{equation}
with $\Lact = (m/((1-f_{\mathrm{air}}^{\max})\rhom))^{1/3}$, the
cube side at the air-fraction saturation. This is one scalar
multiplication per surface point, identical in implementation cost
to the Cauchy formula, and uses a quantity ($\Sab$) that AEGIS
already computes. The Cauchy projection
factor disappears.

A more cautious replacement would present
\cref{eq:replacement} as the operative formula for smooth body parts
and retain the Cauchy formula as a worst-case upper bound applicable
when the cube engulfs a small body feature:
\begin{equation}\label{eq:bound}
  \langle\SAR\rangle^{\mathrm{cube}}
  \;\leq\; \frac{\Sinc\,\Tzero}{\rhom\,L}\,
    \bigl(|k_x|+|k_y|+|k_z|\bigr).
\end{equation}
The bound is loose by up to a factor of $\sqrt{3}$ even in its regime
of best validity, and can be loose by more than a factor of three
otherwise. It is not useful for compliance reporting, only for
order-of-magnitude sanity.

\subsection{What to verify next}

Three numerical checks would close the loop:

\begin{enumerate}[nosep,leftmargin=2em]
  \item\emph{Half-space check.} Implement the 62704 cube-expansion
  algorithm on a voxelized planar half-space at 28\,GHz under
  plane-wave illumination, sweep $\theta\in\{0,15,30,45,60,75\}^\circ$,
  and overlay against \cref{eq:replacement} and the Cauchy formula.
  Expected: agreement with \cref{eq:replacement}, divergence from
  Cauchy as in \cref{tab:compare}.
  \item\emph{Phantom check.} Repeat on the Thelonious or Duke phantom
  with several $\khat$ orientations. The torso, head, and thigh
  should track \cref{eq:replacement}; the ear and finger may diverge.
  \item\emph{Convex-corner check.} Construct a synthetic right-angle
  body corner and verify that both formulas agree there. This
  confirms the Cauchy bound is tight in its claimed regime.
\end{enumerate}

If these checks confirm the analysis, the corollary in
\texttt{paper.tex} should be replaced with \cref{eq:replacement} (or
removed entirely, since for a smooth body the surface-SAR formula
$\SAR(\rr,0)=2\alpha\Sab/\rhom$ already gives the relevant bound,
and the cube quantity differs only by the depth-scale ratio
$\Lact/(2\alpha)^{-1}$).

\subsection{Self-criticism of the standalone derivation}

Two readings of the standalone note's energy argument can be
distinguished, and both are problematic:

\begin{itemize}[leftmargin=1.5em]
  \item Reading A: ``Inward Poynting flux through the cube boundary
  in the thin-skin limit equals $\Sinc A_\perp(\khat)$.'' This
  conflates the cube's geometric shadow with the actual air-facing
  inflow. For a 62704-valid cube the air-facing inflow is much smaller
  than the geometric shadow, because most of the cube boundary is
  immersed in tissue. The conflation is inherited from a setup in
  which the cube is suspended in air with body filling its shadow,
  not embedded in body with surface near one face.
  \item Reading B: ``All photons entering the cube's geometric
  shadow eventually hit the body and are absorbed there, regardless
  of where the cube is.'' This is true for an isolated body, but the
  cube boundary is the wrong control surface. The relevant quantity
  is absorbed power inside the cube, which equals the integral of
  $\Sab$ over the body surface inside the cube footprint, not the
  cube's external shadow.
\end{itemize}

Either way, the cube formula reduces to a surface integral over
$\Sigma\cap\mathcal{C}$, not a projected-area expression on
$\partial\mathcal{C}$. The standalone note's eq.~38 gets this right
in the abstract (``general surface integral $\int_{\Sigma\cap\mathcal{C}}
\Sab\,\diff A$''), then derives a cleaner-looking but inequivalent
expression by replacing the surface integral with a cube-shadow
expression that is only correct in the engulfment limit.


% ======================================================================
\section{Summary}
\label{sec:summary}
% ======================================================================

\begin{enumerate}[leftmargin=2em]
  \item The IEC/IEEE 62704-1 cube is axis-aligned, fixed at $10$\,g
  mass with up to $10\%$ air, with no entirely-air face
  (\cref{sec:standard}). The optimal valid placement on a planar body
  saturates the air constraint, giving cube side
  $\Lact=22.31$\,mm.
  \item Five of the six cube faces are immersed in tissue; only the
  top face (and thin air strips on the four side faces) channels
  fresh wave power into the cube (\cref{sec:audit}).
  \item The absorbed power inside the cube is
  $\Sinc\Tzero\Lact^2\cos\theta_i$, equal to the integral of $\Sab$
  over the cube footprint on the body surface
  (\cref{eq:correct-Pabs}), not $\Sinc\Tzero L^2(|k_x|+|k_y|+|k_z|)$.
  \item The Cauchy projection formula overpredicts $\psSAR$ by
  $50$--$150\%$ at oblique incidence
  ($\theta_i\in[30^\circ,60^\circ]$) and underpredicts by $7\%$ at
  normal incidence (\cref{tab:compare}). It is non-monotone in
  $\theta_i$, contrary to physics.
  \item The Cauchy formula is approximately correct for sharp
  convex corners, axis-aligned incidence, or small body parts
  engulfed by the cube. None of these are the dominant
  compliance scenario (\cref{sec:limits}).
  \item The replacement \cref{eq:replacement} is one scalar
  multiplication per surface point, uses $\Sab$ which AEGIS
  already computes, and matches the standardized $\psSAR$ for
  smooth body parts in the thin-skin regime.
  \item A short numerical experiment with a 62704-style cube
  expansion on a half-space body would settle the matter
  definitively (\cref{sec:implications}).
\end{enumerate}

\begin{thebibliography}{9}

\bibitem{ICNIRP2020}
ICNIRP, ``Guidelines for limiting exposure to electromagnetic fields
(100\,kHz to 300\,GHz),'' \emph{Health Physics}, vol.~118, no.~5,
pp.~483--524, 2020.

\bibitem{62704-1}
IEC/IEEE~62704-1:2017, ``Determining the peak spatial-average
specific absorption rate (SAR) in the human body from wireless
communications devices, 30\,MHz to 6\,GHz -- Part~1: General
requirements for using the finite-difference time-domain (FDTD)
method for SAR calculations,'' 2017.

\bibitem{C953}
IEEE, ``IEEE~C95.3-2021 -- Recommended practice for measurements
and computations of electric, magnetic, and electromagnetic fields
with respect to human exposure to such fields, 0\,Hz to 300\,GHz,''
2021.

\bibitem{CSTdocs}
CST Cable Studio 2013 online help, ``SAR Calculation,''
\url{http://www.mweda.com/cst/cst2013/mergedProjects/CST_CABLE_STUDIO/special_postpr/special_postpr_sar.htm},
accessed May 2026. Documents three averaging methods (IEEE C95.3,
IEEE/IEC 62704-1, CST C95.3) and confirms the $10\%$ air-volume rule
and the no-air-face rule.

\bibitem{psSAR10gNote}
R.~Wydaeghe, ``Peak spatial-average SAR over 10\,g from surface
absorbed power density: a closed-form reduction,''
\texttt{theory/psSAR10g.tex}, April 2026.

\end{thebibliography}

\end{document}
